\chapter{Εισαγωγή}
\label{ch:intro}

\section{Αντικείμενο της εργασίας}
\label{sec:intro-scope}

Η εργασία εξετάζει την επιτάχυνση της επιλογής \tl{top-k}, δηλαδή του
εντοπισμού των \tl{k} μεγαλύτερων στοιχείων ενός συνόλου \tl{n} τιμών, σε τρεις
αρχιτεκτονικές: έναν επεξεργαστή γενικού σκοπού, μια μονάδα γραφικών και μια
μονάδα νευρωνικής επεξεργασίας ενσωματωμένη στην ίδια ψηφίδα με τον
επεξεργαστή.

Για την εκάστοτε αρχιτεκτονική αναπτύσσονται δύο αλγοριθμικές προσεγγίσεις με
ριζικά διαφορετικά χαρακτηριστικά εκτέλεσης. Η περικομμένη διτονική ταξινόμηση
εκτελεί σταθερή ακολουθία συγκρίσεων ανεξάρτητη από τα δεδομένα, ενώ η επιλογή
με τοπικούς σωρούς κατά το σχήμα \tl{map-reduce} προσαρμόζει την εργασία της
στις τιμές της εισόδου. Και οι έξι υλοποιήσεις μετρώνται στον ίδιο παραμετρικό
χώρο, με κοινή υποδομή και σημασιολογία μέτρησης, ώστε οι διαφορές που
προκύπτουν να αποδίδονται στους αλγορίθμους και στο υλικό, όχι στον τρόπο
μέτρησης.

Το βασικό ερώτημα είναι κάτω από ποιες συνθήκες συμφέρει η ανάθεση σε
επιταχυντή, συνυπολογίζοντας τόσο την ταχύτητα της αρχιτεκτονικής, όσο και το
επιπρόσθετο κόστος της διαδικασίας.

\section{Κίνητρο}
\label{sec:intro-motivation}

Οι μονάδες νευρωνικής επεξεργασίας πλέον ενσωματώνονται στην ίδια ψηφίδα με τον
επεξεργαστή των σύγχρονων φορητών υπολογιστών, μοιράζονται την κύρια μνήμη του
και είναι διαθέσιμες για κάθε εφαρμογή που δύναται να τις αξιοποιήσει. Το κατά
πόσο αξίζει να ανατεθεί αυτός ο υπολογισμός στον επιταχυντή αποτελεί πλέον
επίκαιρο ερώτημα για την ανάπτυξη λογισμικού σε τέτοιο υλικό.

Σχεδιάστηκαν, όμως, για συγκεκριμένο φορτίο. Το προγραμματιστικό τους μοντέλο
οργανώνεται κυρίως γύρω από τελεστές νευρωνικών δικτύων και η αρχιτεκτονική
τους γύρω από την παραδοχή ότι κάθε δεδομένο θα χρησιμοποιηθεί πολλές φορές
αφού μεταφερθεί. Τι συμβαίνει όταν το φορτίο δεν έχει αυτή τη μορφή παραμένει
εν πολλοίς ανεξερεύνητο.

Η επιλογή \tl{top-k} είναι κατάλληλο πρόβλημα για να τεθεί αυτό το ερώτημα.
Είναι πρόβλημα απλό στη διατύπωση και διαδεδομένο στην πράξη, καθώς εμφανίζεται
στην αναζήτηση πλησιέστερων γειτόνων, στα συστήματα ανάκτησης και συστάσεων,
στη δειγματοληψία των γλωσσικών μοντέλων και πολλές άλλες εφαρμογές.
Ταυτόχρονα, οι επιταχυντές σχεδιάστηκαν για φορτία με πολλές πράξεις ανά
στοιχείο και εκτενή επαναχρησιμοποίηση δεδομένων, ενώ η επιλογή \tl{top-k} έχει
ελάχιστες πράξεις ανά στοιχείο και καθόλου επαναχρησιμοποίηση.

\section{Συνεισφορά}
\label{sec:intro-contribution}

\paragraph{Υλοποιήσεις.} Κατασκευάστηκαν έξι υλοποιήσεις, δύο αλγόριθμοι σε
τρεις αρχιτεκτονικές, με υποστήριξη πέντε τύπων δεδομένων. Στον επεξεργαστή
αξιοποιούνται οι διανυσματικές επεκτάσεις \tl{AVX-512} και πολλαπλά νήματα. Στη
μονάδα γραφικών η δουλειά οργανώνεται σε πλακίδια με κοινή μνήμη. Η υλοποίηση
της μονάδας νευρωνικής επεξεργασίας γράφτηκε εξ ολοκλήρου στο επίπεδο του
πλαισίου \tl{MLIR-AIE}. Οι υλοποιήσεις επαληθεύτηκαν σε 2700 περιπτώσεις
συνολικά.

\paragraph{Υποδομή μέτρησης.} Κάθε περίπτωση μετριέται σε δύο ρητά διακριτά
πεδία, το αλγοριθμικό και το \tl{end-to-end}. Η κατανάλωση ενέργειας
καταγράφεται με ενιαία μεθοδολογία και στις τρεις αρχιτεκτονικές, ο όγκος των
διακινούμενων δεδομένων μετριέται απευθείας, ενώ η απόδοση κάθε εκτέλεσης
αναλύεται βάσει του μοντέλου οροφής (\tl{Roofline}) που έχει εξαχθεί για το
ίδιο μηχάνημα. Η συγκεκριμένη μεθοδολογία είναι ανεξάρτητη από το πρόβλημα του
\tl{top-k} και μπορεί να εφαρμοστεί σε οποιοδήποτε υπολογιστικό φορτίο
εξετάζεται για ανάθεση σε επιταχυντή.

\paragraph{Ευρήματα.} Η επιλογή \tl{top-k} είναι πρόβλημα κίνησης δεδομένων σε
κάθε αρχιτεκτονική που εξετάστηκε και καμία υλοποίηση δεν περιορίζεται από την
υπολογιστική ικανότητα. Η κατάταξη των δύο αλγορίθμων εξαρτάται από τον λόγο
$n/k$ και παραμένει ίδια σε όλες τις αρχιτεκτονικές. Το πεδίο της μέτρησης
ανατρέπει πλήρως την κατάταξη των αρχιτεκτονικών. Ενεργειακά, η μονάδα
νευρωνικής επεξεργασίας τραβάει τη μικρότερη ισχύ απ' όλες τις υλοποιήσεις, και
παρ' όλα αυτά καταλήγει με ιδιαίτερα αυξημένο ενεργειακό αποτύπωμα, απλώς
επειδή χρειάζεται πολύ περισσότερο χρόνο.

\paragraph{Αναπαραγωγιμότητα.} Ολόκληρη η πειραματική διαδικασία εκτελείται από
δέσμες εντολών που περιλαμβάνονται στο αποθετήριο. Οι συνθήκες μέτρησης
καταγράφονται μαζί με τα αποτελέσματα. Οι λεπτομέρειες δίνονται στο
παράρτημα~\ref{app:reproducibility}.

\section{Δομή της εργασίας}
\label{sec:intro-outline}

Το κεφάλαιο~\ref{ch:background} ορίζει το πρόβλημα, παρουσιάζει τις οικογένειες
αλγορίθμων που το λύνουν, περιγράφει τα εκτελεστικά μοντέλα των τριών
αρχιτεκτονικών και καταλήγει στο κενό που καλείται να καλύψει η εργασία.

Ακολουθεί η μεθοδολογία (κεφάλαιο~\ref{ch:methodology}). Παρουσιάζει τι
μετριέται, με ποια σημασιολογία, με ποια μέτρα σύγκρισης και υπό ποιες
συνθήκες. Αποτελεί προϋπόθεση για την ανάγνωση των αποτελεσμάτων του
κεφαλαίου~\ref{ch:evaluation}.

Έπειτα, ακολουθούν τα κεφάλαια υλοποίησης, ένα ανά αρχιτεκτονική
(κεφάλαια~\ref{ch:cpu}, \ref{ch:gpu}, \ref{ch:npu}), με κοινή εσωτερική δομή
ώστε οι αντίστοιχες ενότητες να συγκρίνονται άμεσα: σχεδιαστικές επιλογές, οι
δύο αλγόριθμοι, διαχείριση μνήμης, υλοποίηση αναφοράς, περιοριστικοί
παράγοντες.

Το κεφάλαιο~\ref{ch:evaluation} παρουσιάζει και ερμηνεύει τα αποτελέσματα,
δηλαδή τον χρόνο εκτέλεσης, τη σύγκριση αρχιτεκτονικών, το αποτύπωμα μνήμης και
την ενεργειακή κατανάλωση, για διαφορετικές παραμέτρους, τύπους δεδομένων και
κατανομές εισόδου.

Το κεφάλαιο~\ref{ch:conclusions} κλείνει με τα συμπεράσματα, τα όρια εντός των
οποίων ισχύουν και πιθανές κατευθύνσεις συνέχειας.

Ακρωνύμια, γλωσσάρι ελληνικών όρων και οδηγίες αναπαραγωγής βρίσκονται στα
παραρτήματα.
